\subsection{模的对称代数}

设 V 是 K-模。V 可唯一地视为交换李代数。于是 V 的包络代数可如下得到：令 T 为 V 的张量代数；令 I 为由张量 \(x \otimes y - y \otimes x\)（\(x \in V, y \in V\)）生成的 T 的双边理想；然后构造代数 S = T/I。

回顾 S 被称为 V 的对称代数（Algebra, Chapter III, § 6），我们简要总结本章所需的性质，其证明都是直接的。令 \(T^{n}\) 为 T 中 n 阶齐次张量的集合。则 \(\mathrm{I} = (\mathrm{I} \cap \mathrm{T}^{2}) + (\mathrm{I} \cap \mathrm{T}^{3}) + \cdots\)，从而 S 是各 \(S^{n}\)（它们是 \(T^{n}\) 的典范像）的直和。\(S^{n}\) 的元素称为 n 次齐次元素。\(S^{0} = K.1\)，\(S^{1}\) 与 V 相同，且 \(S^{n}S^{p} \subset S^{n+p}\)。代数 S 由 1 和 \(S^{1} = V\) 生成。显然 \(S^{1}\) 中任意两个元素都可交换，因而 S 是交换的。若 V 是以 \((x_{\lambda})_{\lambda \in \Lambda}\) 为基的自由 K-模，则从多项式代数 \(K[X_{\lambda}]_{\lambda \in \Lambda}\) 到 S 的典范满同态 f（将 1 映到 1，并将 \(X_{\lambda}\) 映到 \(x_{\lambda}\)，对所有 \(\lambda \in \Lambda\)）是同构：因为由 S 的泛性质（第 1 号，命题 1），存在从 S 到 \(K[X_{\lambda}]_{\lambda \in \Lambda}\) 的同态 g，将 1 映到 1，并将 \(x_{\lambda}\) 映到 \(X_{\lambda}\)，对所有 \(\lambda \in \Lambda\)，且 f 与 g 互为逆同态。

令 \(S'^{n} \subset T^n\) 为阶数 \(n\) 的齐次对称张量的集合（Algebra, Chapter III, § 5, no. 1, Definition 2）。若 \(K\) 是特征为 0 的域，则 \(S'^{n}\) 与 \(I \cap T^n\) 在 \(T^n\) 中互为补空间。事实上，设 \((x_{\lambda})_{\lambda \in \Lambda}\) 是 \(V\) 的一个基。给 \(\Lambda\) 赋予全序（Set Theory, Chapter III, § 2, no. 3, Theorem 1）。令 \(\Lambda_n\) 为由 \(n\) 个 \(\Lambda\) 中元素组成的递增序列的集合。对 \(M = (\lambda_1, \ldots, \lambda_n) \in \Lambda_n\)，令

\[
y _ {\mathrm{M}} = \frac {1}{n !} \sum_ {\sigma \in \mathfrak {S} _ {n}} x _ {\lambda_ {\sigma (1)}} \otimes \dots \otimes x _ {\lambda_ {\sigma (n)}}.
\]

\(y_{\mathbf{M}}\) 对于 \(\mathbf{M} \in \Lambda_n\) 构成向量 K-空间 \(\mathbf{S}'^n\) 的一组生成元。根据上段，它们在 \(\mathbf{S}^n\) 中的典范像构成 \(\mathbf{S}^n\) 的一个基。因此 \((y_{\mathbf{M}})_{\mathbf{M} \in \Lambda_n}\) 是 \(\mathbf{I} \cap \mathbf{T}^n\) 在 \(\mathbf{T}^n\) 中的一个补子空间的基（Algebra, Chapter II, § 1, no. 6, Proposition 4），这就证明了断言。

因此，当 K 是特征为 0 的域时，\(S'^{n}\) 上的典范映射 \(T^{n} \rightarrow S^{n}\) 的限制是从空间 \(S'^{n}\) 到空间 \(S^{n}\) 的同构，因而有逆同构。对每个 n 得到的这些逆同构定义了从空间 S 到对称张量空间 \(S' = \sum_{n \geqslant 0} S'^{n}\) 的典范同构。

